\documentclass[11pt]{article}
\usepackage[margin=0.8in]{geometry}
\usepackage[T1]{fontenc}
\usepackage{amsmath,amssymb,booktabs,xurl,hyperref}
\setlength{\emergencystretch}{3em}
\title{Root Causes in the Eleven Former CLF-Interrupted Scenarios}
\author{Pan Dynamics Collision Avoidance Research}
\date{October 1, 2026}
\begin{document}
\maketitle
\section{Finding and scope}
Removing the solver-status admission gate fixed a numerical interruption,
not the controller's avoidance and recovery behavior. This diagnosis covers
all eleven formerly CLF-interrupted cases from the direct two-stage campaign.
Three later stop at an infeasible PCBF problem. Eight continue for 80 s without
nominal recovery; two of those eight also collide.

The common structural issue is how an affine prediction becomes the next
linearization reference under fixed hard correction boxes. A dynamically
inconsistent reference can make the next subproblem infeasible. Its shifted
future inputs can also exclude the current recovery action even when the
CLF objective is being minimized correctly. Positive PCBF relaxation then
allows collision-containing local plans to be issued. Nonlinear prediction
error is an additional measured issue, not a substitute explanation for a
plan that already predicts a collision.

No production controller, configuration, online check, retry, fallback or
mode was added. All interventions below are offline diagnostic ablations;
none is presented as a validated replacement algorithm.

\section{Methods and coverage}
The controller under study is commit
\path{0466fd95800e80d53434cf90b8021da5989932b0}; the complete campaign is
reported by commit \path{efc5e1d33d4dfa76c6d59c096105babdc49a85d5} in
\path{report/DIRECT_TWO_STAGE_SCENARIOS_20261001.tex}.
Conditions remain exact state and known target motion, no road constraints,
6-mm collision buffer, 50-ms holds, horizons of 68/76 holds at 8/15 m/s,
and one PCBF solve followed by one CLF solve. No random draws are used.

The three later failure problems are captured at their actual failing calls.
For the eight 80-s cases, the final saved physical state and previous affine
prediction reconstruct the next hypothetical call at 80.00 s. These are
one-call diagnostics, not extra executed campaign frames. Earlier contexts
are reconstructed by replaying the controller against saved actual states.
All 201 replayed inputs for 8-m/s acceleratingHeadOn, 22 for 15-m/s headOn,
and 21 for 15-m/s acceleratingHeadOn match the original inputs exactly.
Additional onset replays also match exactly.

Instrumentation copies the unchanged formulation, labels constraint rows,
and captures matrices before solving. Offline coneprog uses 400 iterations,
constraint tolerance $10^{-8}$, optimality tolerance $10^{-7}$ and a 5-s soft
budget. Linear feasibility and Farkas probes use linprog tolerances $10^{-9}$.
Where stated, ODE45 evaluates a held input or complete planned sequence with
relative/absolute tolerances $10^{-11}/10^{-12}$. These evaluations do not
run in the controller. No new runtime benchmark is inferred from diagnostic
calls. Source/native identity was verified against the frozen campaign copy.

\begin{center}\small
\begin{tabular}{r l l}
\toprule
Speed (m/s) & Scenario & Observed outcome / diagnosed mechanism\\
\midrule
8 & headOn & PCBF failure at 2.35 s; dynamics/trust conflict\\
8 & acceleratingHeadOn & 80 s unrecovered; recovery input restricted\\
8 & brakingLead & PCBF failure at 11.40 s; dynamics/trust conflict\\
8 & turningCrossing & 80 s unrecovered; recovery input restricted\\
8 & curvedCrossing & 80 s unrecovered; recovery input restricted\\
15 & headOn & Collision; also unrecovered at 80 s\\
15 & acceleratingHeadOn & Collision; also unrecovered at 80 s\\
15 & brakingLead & PCBF failure at 23.95 s; dynamics/trust conflict\\
15 & crossing & 80 s unrecovered; recovery input restricted\\
15 & turningCrossing & 80 s unrecovered; recovery input restricted\\
15 & curvedCrossing & 80 s unrecovered; recovery input restricted\\
\bottomrule
\end{tabular}\end{center}
The other three original PCBF-interrupted scenarios are excluded from this
new set of eleven; their causes are documented in the preceding reports.

\section{Eight unrecovered cases: the optimizer cannot choose the return input}
\subsection{The restriction is centered on a future plan, not the actuator}
The implementation shifts the prior affine states and inputs without
regenerating the state sequence from the new physical state. It then imposes
\[
 |\delta_i-\bar\delta_i|\le0.075\ \mathrm{rad},\qquad
 |b_i-\bar b_i|\le0.125.
\]
Here $\bar u_0$ is the previous plan's \emph{second} input. These are hard
local approximation restrictions, not the physical actuator limits
($|\delta|\le40^\circ$, $|b|\lesssim1$). The tested configuration has no finite
actuator slew limit, so these cannot be justified as a measured slew bound.

At 80 s in 8-m/s acceleratingHeadOn, the ego is 354.07 m above its path with
heading error $0.7582$ rad. The shifted steering anchor is $0.0752319$ rad,
so the first steering must lie in
\[
 [0.0002319,\ 0.1502319]\ \mathrm{rad}.
\]
The optimizer requests approximately the lower bound, $0.0002339$ rad.
It cannot choose a negative steering input to turn back. The situation was
already established at 10 s: the allowable lower endpoint and returned
steering are approximately $+0.0013436$ rad, while lateral error is 35.30 m.
The problem is not that the CLF solver elected to stop improving too early.
It has reached the useful edge of the allowed input set.

\begin{center}\small
\begin{tabular}{r l r r r r}
\toprule
Speed & Scenario & $e_y$ (m) & Anchor $\bar\delta$ & Allowed $\delta$ interval & Returned $\delta$\\
\midrule
8 & acceleratingHeadOn & 354.07 & 0.07523 & $[0.00023,0.15023]$ & 0.00023\\
8 & turningCrossing & -296.94 & -0.07455 & $[-0.14955,0.00045]$ & 0.00045\\
8 & curvedCrossing & -224.23 & -0.07315 & $[-0.14815,0.00185]$ & 0.00185\\
15 & headOn & 725.12 & 0.07970 & $[0.00470,0.15470]$ & 0.00470\\
15 & acceleratingHeadOn & 727.14 & 0.08218 & $[0.00718,0.15718]$ & 0.00727\\
15 & crossing & -628.85 & -0.08686 & $[-0.16186,-0.01186]$ & -0.01186\\
15 & turningCrossing & -728.69 & -0.08280 & $[-0.15780,-0.00780]$ & -0.00793\\
15 & curvedCrossing & -373.19 & -0.08168 & $[-0.15668,-0.00668]$ & -0.00672\\
\bottomrule
\end{tabular}\end{center}
Steering values are in radians. The two collision cases' late states belong
to the simulation without impact dynamics, not a physical post-impact model.

\subsection{Independent optimization and causal interventions}
Eliminating the first affine transition expresses the CLF objective as a
quadratic in just two first-input increments,
\[
 V_1(d)=\alpha\|Gd+c\|^2,\qquad
 \rho=\max\{0,V_1(d)-0.99V_0\}.
\]
Minimizing this quadratic exactly over its two-dimensional box, with all
future constraints removed, gives a lower bound for the complete problem.
For all eight late cases the exact minimizer is at both input-box boundaries
(the steering direction shown above and the lower braking-ratio endpoint).
The complete conic solutions nearly attain that lower bound: absolute CLF
slack gaps are 0.0059--3.3527 against slack values of about
$5.10\times10^4$--$3.58\times10^5$. All eight reconstructed calls return
positive solver flags for both stages.

Removing the terminal cone, removing future trust restrictions while keeping
the first-input box, or refreshing only the first-state tangent does not
restore meaningful first-input steering. This rules out those mechanisms as
the immediate cause of the late weak steering. Recentring just the first
input box on the previous \emph{applied} input permits corrective steering in
all eight diagnostic calls and reduces the actual one-step CLF growth.
That intervention alone is not evidence of complete recovery.

There is a stronger earlier counterfactual. At 2.00 s in 8-m/s
acceleratingHeadOn, $V_0=4124.991$. The original problem has essentially zero
PCBF optimum, CLF slack 68.606, predicted $\Delta V=+27.356$, and actual
$\Delta V=+20.138$. Relaxing only the first-input trust box to the unchanged
physical limits, while retaining all other horizon constraints, gives CLF
slack about $1.75\times10^{-4}$, predicted $\Delta V=-45.664$, and actual
$\Delta V=-64.326$. The 1\% requirement is $\Delta V\le-41.250$.
The actual sampled clearance over that hold remains at least 1.214 m.
Thus the original input restriction excludes an available dissipating action
at this state. This is a one-hold causal result, not a full alternative rollout.

\subsection{Why the reference keeps creating the restriction}
The second stage minimizes only a one-step CLF slack. The future controls
have no nominal-recovery objective of their own. They participate in
feasibility and terminal admission, but when multiple tails have the same
first-step objective, the objective does not select the one that provides a
useful next linearization input. The free-pose terminal core constrains
intrinsic motion, not return to the given path. Shifting this underdetermined
tail and using it as the center of a hard input box promotes a numerical
future choice into a restriction on the next real action.

A positive slack also allows $V$ to increase. Minimizing it cannot establish
convergence unless the admissible input set allows dissipation. The quadratic
matrix was derived by LQR around a nominal trim; the code does not establish
a global CLF property for hundreds of metres of lateral error and large
heading errors. Even removing all trust bounds in the late affine problems
does not make the prescribed 1\% contraction feasible. Moreover, the large
steps then proposed can perform worse in the actual nonlinear model.
Therefore wholesale trust-bound deletion is not justified as a fix. The
measured early exclusion and late reference trapping are concrete failures;
a global convergence theorem is not claimed in either direction.

\section{Three later PCBF failures: artificial infeasibility from shifted states}
The local affine equation is
\[
 d x_{i+1}-A_i d x_i-B_i d u_i
 =F_h(\bar x_i,\bar u_i)-\bar x_{i+1}.
\]
Because $\bar x$ is a shifted affine prediction rather than a rollout,
the right-hand side can be large. The state corrections are restricted to
$\pm(2.5,2.5,0.25,2.5,1.5,0.75)$ in their respective units, and the input
corrections to $\pm(0.075,0.125)$. Repairing all dynamics defects simultaneously
can then be impossible.

\begin{center}\small
\begin{tabular}{r l r r l}
\toprule
Speed & Scenario & Failure time & Farkas value & Largest one-step defect\\
\midrule
8 & headOn & 2.35 s & -0.1698601 & yaw rate: 0.777009 rad/s, stage 16\\
8 & brakingLead & 11.40 s & -0.0503176 & yaw rate: 0.231437 rad/s, stage 1\\
15 & brakingLead & 23.95 s & -0.0344165 & lateral speed: 1.029350 m/s, stage 19\\
\bottomrule
\end{tabular}\end{center}
These defect magnitudes are differences between successor states, not rates
of change of the listed quantities. In all three cases:
\begin{itemize}
\item Removing collision rows, terminal conditions, or the CLF cone individually
  does not restore feasibility.
\item Keeping \emph{only} dynamics and state/input trust boxes is already
  infeasible as a linear program, without any safety, road or terminal rows.
\item Zeroing dynamics defects in a diagnostic problem restores linear
  feasibility; zeroing only the initial-state mismatch does not.
\item Removing input trust bounds while preserving physical actuator/state
  limits restores conic feasibility, with PCBF optima below $2.2\times10^{-7}$.
\end{itemize}
Farkas witnesses for the original linear subsystems have nonnegative inequality
multipliers, negative combined right-hand sides as tabulated, and equality
residuals at most $2.23\times10^{-16}$. They distinguish mathematical
inconsistency from a slow or unsuccessful numerical search.

A fresh flow seed at the same physical state restores both solver stages in
the two 8-m/s cases. It does not restore the 15-m/s brakingLead case. Seed
reset is therefore not a universal cure; the conflict between model defects
and hard neighborhoods is the established cause in these three captured
problems. None of these results establishes physical unavoidability.

\section{Two collisions: lost zero-slack feasibility followed by relaxed execution}
The 15-m/s headOn and acceleratingHeadOn cases first overlap at approximately
1.09833 s and 1.05000 s, respectively, in the previous dense sampled audit.
The corresponding control calls are at 1.05 s and 1.00 s.

\subsection{The local model loses zero-slack feasibility before contact}
The first materially positive PCBF optima appear at 0.60 s for headOn and
0.65 s for acceleratingHeadOn:
\begin{center}\small
\begin{tabular}{l r r r}
\toprule
Case / call time & Original PCBF optimum & Reroll same input anchor & New flow anchor\\
\midrule
headOn / 0.60 s & $2.3843\times10^{-3}$ & $5.0832\times10^{-8}$ & $2.6636\times10^{-8}$\\
acceleratingHeadOn / 0.65 s & $8.3791\times10^{-4}$ & $1.7731\times10^{-8}$ & $2.2923\times10^{-8}$\\
\bottomrule
\end{tabular}\end{center}
Rerolling means retaining the same shifted input sequence, integrating it from
the current actual state, and rebuilding dynamics and geometry around those
states. It is a diagnostic reconstruction, not a production validation step.
Forcing zero slack in the original problems is infeasible. Removing terminal
conditions, hard tail collision rows, or state trust bounds individually does
not remove their positive optima. Removing input trust bounds nearly removes
them; changing the reference removes them without changing physical limits.
These facts implicate the local reference, its geometry, and restricted input
corrections rather than proving no avoidance trajectory exists physically.

\subsection{Even near-zero affine slack is not equivalent to the nonlinear plan}
At 0.55 s in headOn, the original PCBF optimum is
$1.08\times10^{-8}$ and the returned affine trajectory has minimum node
clearance 9.315 mm. Holding its planned input sequence in an independent
nonlinear rollout produces strict overlap: the minimum separating-axis gap is $-4.172$ mm
at absolute time 1.10 s. Maximum position disagreement over that
3.8-s rollout is 0.368 m. At 0.60 s in acceleratingHeadOn, the corresponding
affine minimum node clearance is 11.040 mm, while the nonlinear rollout's
start/midpoint/end samples reach 5.896 mm, below the configured 6-mm margin.
These are full planned-sequence counterfactuals, not the actual future
receding-horizon trajectory. They demonstrate the missing implication from
affine feasibility to the nonlinear continuation.

\subsection{The actual collision inputs already come from collision predictions}
At the eventual collision-producing calls, the optimum slack sums are
0.023411 and 0.006588. Both returned affine endpoint rectangles already
overlap the target; the actual ODE45 endpoint rectangles also overlap.
The first input's endpoint is constrained at the next stage, whose safety
slack may be positive. A nearly zero first-stage slack is therefore not a
safe-next-endpoint test.

The first-step affine position errors at these two calls are only about
0.044 mm and 0.0094 mm in Euclidean norm. Thus these particular impacts
cannot be described as a large unpredicted first-step deviation. The optimizer
already accepts a relaxed collision-containing prediction. Safety slack sums
are aggregate residuals, not a required clearance in metres or a direct
penetration depth. Once positive slack is the best local result, its numerical
existence alone does not make collision avoidance successful.

\section{Design implications and limits}
The evidence supports three connected design priorities within the existing
two-stage architecture: generate a dynamically consistent and useful
linearization reference; avoid turning underdetermined future controls into
hard exclusions of the next recovery action; and distinguish minimization of
a collision relaxation from actually achieving avoidance. The CLF also needs
an admissible recovery mechanism beyond an unsupported global use of a
local one-step quadratic. These are formulation issues; longer simulation,
more solver time, gate removal, road removal, or changing the terminal cone
alone does not address the observed mechanisms.

No modified controller was validated here. In particular, directly removing
all trust bounds can produce inaccurate large-step predictions, and changing
a seed at one call does not prove an entire scenario will succeed. No extra
online nonlinear checking or backup architecture is required by this diagnosis.

\section{Artifacts and verification}
The accompanying directory
\path{report/ELEVEN_SCENARIO_FAILURE_CAUSES_20261001/}
contains captured-problem statistics, eight-case recovery bounds,
27 later-failure ablations, three Farkas witnesses, collision-onset and
pre-impact counterfactuals, exact replay comparisons, diagnostic scripts and an instrumentation generator,
actual console exports and source/raw hashes. Console exports omit terminal
backspace characters and trailing whitespace; original logs remain in the
cache with recorded hashes. Large MAT captures remain in
\path{/home/zai/.cache/collisionAvoidance/eleven-failure-causes-20261001/}.
The generated instrumentation is research evidence and is not on the production
controller path. MATLAB scripts executed as recorded in the logs; this report
was compiled and text extraction checked. No production unit suite was rerun
for the unchanged controller. All claims about recovery are limited to the
observed trajectories and the explicitly identified diagnostic calls.
\end{document}
